%!TEX root = ../../thesis.tex

\section{Fuzzing}\label{sec:fuzzing}

% Drafting outline: bullet points to be fleshed out into prose.
% \todo{cite: ...} marks claims whose canonical reference is not yet in references/*.bib.

\subsection{Overview and Motivation}\label{sec:fuzzing-overview}

% --- Verbatim excerpt from the proposal (own text; citation keys remapped to references/*.bib) ---
Fuzz testing, or \textit{fuzzing} for short, is a widely deployed method for automatic
software testing and vulnerability discovery that makes use of a diverse set of
tools~\cite{manes_art_2019}.
Its goal is to find inputs to a System Under Test (SUT) that trigger a certain
behavior, often a crash that might then lead to the discovery of bugs or security
vulnerabilities.
% --- End excerpt ---

The technique was first introduced in the early 1990s as a way of testing the
reliability of UNIX utilities~\cite{miller_empirical_1990}.
The authors were inspired by spurious characters being introduced to a dial-up
line causing crashes of common programs on a remote workstation.
This led to the creation of a simple command-line utility \emph{fuzz} that produced
random input strings that were then sent to various utilities in the hope of causing
a crash in the form of a core dump or a hang.
They managed to find inputs that caused between 24\% and 33\% of the tested
utility programs to crash or hang.
Analyzing these inputs identified bugs, which could then be fixed, thereby
contributing to the reliability of the system.

The same approach is still in use today to find security vulnerabilities in
input-processing code, especially memory-safety bugs in C/C++.
For example, Google's OSS-Fuzz
\enquote{aims to make common open source software more secure and stable by combining
modern fuzzing techniques with scalable, distributed execution}~\cite{arya_oss-fuzz_nodate}.
It supports a variety of target languages and is therefore widely applicable.
OSS-Fuzz has helped to identify over \num{13000} security vulnerabilities and \num{50000} bugs across
\num{1000} projects, proving the importance of modern fuzzing techniques for ensuring reliability and security.

\subsection{Core Terminology}\label{sec:fuzzing-definitions}

We will introduce some general definitions that mainly follow \textcite{manes_art_2019}.

\begin{enumerate}
  \item \emph{System Under Test (SUT):}\index{system under test}
    The system under test is an executable software artifact like a command-line utility
    or library together with an interface that consumes input. In our case it denotes
    a binary executed through the MOGI emulator (\cref{chap:implementation}).

  \item \emph{Security Policy:}\index{security policy}
    A security policy defines which observable behaviors of the SUT are acceptable.
    Each execution either satisfies or violates a given security policy. A common
    security policy deems all executions that result in crashes like segmentation
    faults to be violations.

  \item \emph{Bug Oracle:}\index{bug oracle}
    A bug oracle is a program that decides whether a given input to the SUT violates
    a given security policy. A widely used bug oracle might deem an execution
    a violation if the SUT crashes or hangs.

  \item \emph{Fuzzing:}\index{fuzzing}
    Fuzzing denotes the process of feeding the SUT with inputs sampled from an input space.
    This input space may include expected and valid inputs as well as invalid, unexpected
    ones, which are sampled in the hope of violating a given security policy.

  \item \emph{Fuzz Testing:}\index{fuzz testing}
    Fuzz testing denotes the software testing technique that makes use of fuzzing.

  \item \emph{Fuzzer:}\index{fuzzer}
    A fuzzer is a program that performs fuzz testing on a SUT.

  \item \emph{Fuzz Campaign:}\index{fuzz campaign}
    A fuzz campaign is a specific execution of a fuzzer on a SUT with a specific
    security policy. Its goal is to find bugs in the SUT that violate this policy.

  \item \emph{Test Case:}\index{test case}
    A test case denotes a concrete input that the fuzzer feeds to the SUT during
    a fuzzing campaign. The resulting execution is then judged by the bug oracle
    and the test case is recorded if the bug oracle deems it a violation of the
    security policy.
    In our case, depending on the action model, this can either be a single byte or a 
    complete byte sequence.

  \item \emph{Seed:}\index{seed}
    A seed is a commonly well-structured input to the SUT that the fuzzer uses
    as the basis for deriving new test cases by modifying it.

  \item \emph{Corpus:}
    The corpus\index{corpus} is the collection of seeds the fuzzer maintains. Coverage-guided
    fuzzers evolve it over the course of a campaign by retaining test cases that
    exhibit new behavior. We write \emph{initial corpus} for the collection the
    campaign starts from, which is usually supplied by the user.\index{corpus!initial}

  \item \emph{Fuzz Configuration:}\index{fuzz configuration}
    A fuzz configuration comprises the parameter values that control the fuzzing
    algorithm.
\end{enumerate}

\subsection{Taxonomy}\label{sec:fuzzing-taxonomy}

Fuzzers can be roughly divided into four classes, defined by the amount of
information from the execution of the SUT available to them~\cite{liang_fuzzing_2018}.

\subsubsection{Black-box fuzzing}

Black-box fuzzers\index{fuzzing!black-box} have no access to the internal behavior of a SUT during execution.
They can only make use of the bug oracle to decide whether a given input violates
the security policy or not.
The core advantage of this approach is the wide applicability since no instrumentation
is required.
This lack of potential instrumentation overhead also benefits execution throughput
as well as timing and reproducibility for targets like network interfaces.
The disadvantage is the lack of information about SUT behavior.
This makes it difficult for fuzzers to generate high-coverage inputs.

\subsubsection{White-box fuzzing}

In contrast to pure black-box fuzzers, white-box fuzzers\index{fuzzing!white-box} have full access to
the internal logic of the SUT as well as usually any original source code or
documentation which can greatly aid the fuzzing process.
In practice, this allows an approach known as \emph{symbolic execution}\index{symbolic execution}: instead
of relying on actual execution traces, a white-box fuzzer can gather symbolic
constraints at branching points and use constraint solving to find inputs that exercise
a specific path.
While this makes it easy to find test cases of interest in theory, this approach scales poorly
to large SUTs due to the path explosion problem.\index{path explosion}
A popular white-box fuzzer is SAGE~\cite{godefroid_sage_2012}.

\subsubsection{Greybox fuzzing}

Since full information is often not available but observing certain parts
of the internal behavior of the SUT might still be possible, \textit{greybox}
fuzzing\index{fuzzing!greybox} represents an important middle ground between the previous two extremes.
By making use of runtime instrumentation information like code coverage, the fuzzer can
better judge the behavior of the SUT for a given input and direct its exploration efforts towards promising regions.

\subsubsection{Hybrid approaches}

There also exists a variety of hybrid approaches combining different fuzzing methodologies.
A prominent example is \emph{Driller}\index{Driller}, which alternates pure greybox fuzzing with
\emph{concolic} execution\index{concolic execution} (a combination of concrete and symbolic execution) to overcome difficult branch
conditions like CRC checks that have historically presented a major obstacle for fuzzers~\cite{stephens_driller_2016}.

\subsection{Input generation}\label{sec:fuzzing-input-generation}

% Outline follows Manès et al. §5 (Input Generation); the § pointers below refer to that paper.

The content of a test case directly controls whether a bug is triggered or not, making
the way content is generated one of the most influential design decisions of a fuzzer~\cite{manes_art_2019}.
Input generation can be categorized into either generation- or mutation-based.

\subsubsection{Generation-based input generation}

Generation-based input generation\index{input generation!generation-based} makes use of a model of the input expected by
the SUT, often represented through a grammar specification.
These models can be \emph{predefined}, \emph{inferred}, or \emph{encoder-based}.

Predefined models are specified by the user.
They are often used by network fuzzers consuming protocol
specifications~\cite{pereyda_boofuzz_nodate} or kernel fuzzers consuming system
call interfaces.

Inferred models do not rely on an explicitly specified model but rather infer it through some automatic means.
Inference can either take place as a preprocessing step or during a fuzz campaign.
TestMiner~\cite{toffola_saying_2017} searches for literals inside the SUT and derives suitable inputs.
PULSAR~\cite{thuraisingham_pulsar_2015} infers a network protocol from a set of captured network packets generated by the SUT.

Encoder-based models interpret the SUT as a decoder for inputs generated by an existing encoder.
An example for an encoder-based fuzzer is MutaGen~\cite{kargen_turning_2015}.

\subsubsection{Mutation-based input generation}

Generating purely random inputs is not an effective strategy beyond the
shallowest layers of a SUT~\cite{manes_art_2019,godefroid_sage_2012}.
In general, a SUT contains checks that reject invalid inputs, making it
difficult to reach code paths that violate a given security policy~\cite{manes_art_2019}.
This motivates the use of an \emph{initial corpus}, a set of well-structured inputs
to the SUT like valid network packets or files.
The seeds inside the corpus are then successively mutated\index{input generation!mutation-based} with the goal
of triggering abnormal behavior.
Mutation operators can be divided into four sets:

\begin{enumerate}
  \item \emph{Bit-flipping:}
    Bit flipping\index{mutation!bit-flipping} is governed by a \emph{mutation ratio}\index{mutation ratio}, the fraction of bits to
    flip per execution.

  \item \emph{Arithmetic mutation:}
    Arithmetic mutation\index{mutation!arithmetic} reinterprets byte sequences as integers and perturbs them.

  \item \emph{Block-based mutation:}
    Block-based mutation\index{mutation!block-based} inserts, deletes, replaces, permutes or resizes byte blocks,
    and splices blocks across seeds.

  \item \emph{Dictionary-based mutation:}
    Dictionary-based mutation\index{mutation!dictionary-based} substitutes values with other values of high semantic
    weight. These usually include numbers like $0$ or $-1$, or format strings for
    string substitution.
\end{enumerate}

\subsection{Coverage-Guided Greybox Fuzzing}\label{sec:fuzzing-cgf}

This work will focus primarily on coverage-guided greybox fuzzing,\index{fuzzing!coverage-guided greybox} \emph{CGF} for short.
It is the dominant design among modern general-purpose fuzzers, popularized by
AFL~\cite{zalewski2014afl} and its successor AFL++~\cite{fioraldi_afl_nodate}.
We will outline the general approach following \textcite{bohme_coverage-based_2016}.

\subsubsection{Overview}

A coverage-guided greybox fuzzer in general uses mutational input generation
(\cref{sec:fuzzing-input-generation}).
For this, it keeps a corpus $T$ as part of its runtime state.
This corpus is initialized from an initial corpus $S$.
The fuzzer then selects a seed $t\in T$ and derives a new test case $t^\prime$.
If the fuzzer deems the execution of the SUT under $t^\prime$ to be interesting, $t^\prime$
is retained in the corpus $T$ and is thereafter available as a seed itself.

\begin{algorithm}[H]
  \caption{Coverage-guided greybox fuzzing, following \textcite{bohme_coverage-based_2016}.}
  \label{alg:cgf}
  \begin{algorithmic}[1]
    \Require Initial corpus $S$, system under test $P$
    \Ensure Crashing test cases $T_\times$, corpus $T$
    \State $T \gets S$
    \State $T_\times \gets \emptyset$
    \If{$T = \emptyset$}
      \State $T \gets \{\,\text{empty input}\,\}$
    \EndIf
    \Repeat
      \State $t \gets \Call{ChooseNext}{T}$
        \Comment{seed scheduling}
      \State $p \gets \Call{AssignEnergy}{t}$
        \Comment{power schedule}
      \For{$i \gets 1$ \textbf{to} $p$}
        \State $t^\prime \gets \Call{MutateInput}{t}$
        \If{$t^\prime$ crashes $P$}
          \State $T_\times \gets T_\times \cup \{t^\prime\}$
            \Comment{bug oracle}
        \ElsIf{$\Call{IsInteresting}{t^\prime}$}
          \State $T \gets T \cup \{t^\prime\}$
            \Comment{coverage check}
        \EndIf
      \EndFor
    \Until{timeout reached or campaign aborted}
    \State \Return $T_\times$, $T$
  \end{algorithmic}
\end{algorithm}

\Cref{alg:cgf} states the CGF loop.
The fuzzer must solve three main problems: selecting an appropriate
seed (\textsc{ChooseNext}, \textsc{AssignEnergy}), applying appropriate mutations (\textsc{MutateInput}), and deciding whether to retain the
mutated input or not based on SUT behavior (\textsc{IsInteresting}).
In addition, \textsc{AssignEnergy} decides how much of the execution budget the selected seed receives.

Two independent decisions are made after every execution.
The bug oracle decides whether the execution is a
\emph{violation} and therefore a reportable finding (line 12).
Separately, coverage information
is used to decide whether the execution is \enquote{interesting} and therefore worth keeping as a starting point for
further mutation (line 14).

The fuzzing loop can be interpreted as an evolutionary algorithm~\cite{she_neuzz_2019,manes_art_2019}.
The corpus $T$ is a population, the mutation operators of
\cref{sec:fuzzing-input-generation} are the genetic operators, and coverage novelty is the
fitness function deciding which offspring survive into the next generation.
Importantly, the fuzzer does not form an explicit model of what constitutes a good input.
It only observes which test cases reach a location no previous test cases reached before.
This is a very weak learning signal, and a large part of the research surveyed in
\cref{sec:fuzzing-limitations} and \cref{chap:related-work} can be read as an attempt to
extract more information from the same executions.

\subsubsection{Test case selection}

\textsc{ChooseNext} decides which seed $t \in T$ is mutated next, and
\textsc{AssignEnergy} decides how many mutants are derived from it before the fuzzer moves
on.
The two together determine how a finite execution budget is distributed over the test case
set.
They are commonly referred to as \emph{seed scheduling}\index{seed scheduling} and the \emph{power
schedule}\index{power schedule}~\cite{manes_art_2019,bohme_coverage-based_2016}.

An optimal seed and power schedule would choose seeds in such a way as to maximize coverage
after they are mutated.
This coverage is not known to the fuzzer unless it executes the SUT with a mutated seed,
that is, only once the budget has already been spent on that seed.
Scheduling is therefore subject to what \textcite{manes_art_2019} call the
\emph{exploration versus exploitation conflict}\index{exploration--exploitation tradeoff}: time can be spent either on gathering more
accurate information about the available choices in order to inform later decisions
(\emph{explore}), or on fuzzing the ones currently believed to lead to the most favorable
outcome (\emph{exploit}).

\subsubsection{Test case mutation}

\textsc{MutateInput} derives a new test case $t^\prime$ from the selected test case $t$ by
applying one or more mutation operators.
Each mutation requires a location in the input to mutate as well as an actual mutation operation.

The main problem is the size of the mutation space.
For an input of length $l$ and a mutation set $M$, the space of all possible mutations has size $\mathcal{O}(l\lvert M\rvert)$.
This presents a challenge for long inputs as the mutation space is very large compared to the number of useful mutations.
Branch conditions are usually affected by only a few bytes of a given input, making it hard to select appropriate mutations.
Some fuzzers try to derive that dependency through data-flow
analysis~\cite{gan_greyone_nodate} and through neural surrogates that
approximate the
byte-to-coverage map and follow its gradient~\cite{she_neuzz_2019,she_mtfuzz_2020}.

Another problem lies in the existence of magic bytes\index{magic bytes} or checksums\index{checksum}, which are hard to satisfy by random mutations.
For example, a four-byte comparison is satisfied with probability $2^{-32}$ per attempt, making it very unlikely for a single mutation to overcome it.
This motivates a hybrid approach that makes use of concolic execution~\cite{stephens_driller_2016}, combining mutations with highly targeted constraint solving.

Choosing appropriate mutations is a scheduling problem that most conventional fuzzers try to solve through a fixed, hand-tuned distribution.
This has the major drawback that useful mutations are highly target-dependent.
Some fuzzers like MOpt~\cite{lyu_mopt_nodate} therefore search for a better distribution online using particle swarm
optimization, and Rainfuzz~\cite{binosi_rainfuzz_2023} learns a heat map over byte positions with
a policy-gradient method.

\subsubsection{Test case retention}

\textsc{IsInteresting} decides which mutated test cases to keep.
Its decision is based on coverage information.
Coverage information is recorded by a lightweight program instrumentation framework during execution.
The instrumentation can collect coverage at different granularities:

\begin{enumerate}
  \item \emph{Basic-block coverage}\index{coverage!basic-block}
    records which blocks were executed.

  \item \emph{Edge} or \emph{branch coverage}\index{coverage!edge} records which transitions $(A,B)$
    between blocks were taken and is strictly more informative.

  \item \emph{Path coverage}\index{coverage!path} tracks a sequence of block hits $(A,B,\dots)$.
    While this contains the most useful information, execution paths
    can grow large and therefore slow down the fuzzer.
\end{enumerate}

Orthogonal to the granularity is the question of how precisely a hit is counted.
One bit per coverage element is the cheapest option, but it cannot distinguish how often a block was hit.
Most coverage-guided fuzzers therefore keep a hit count per element and quantize it into
buckets of logarithmically scaling size,\index{coverage!hit count bucket} so that a change in hit counts registers as new behavior once a new bucket is entered.

To decide whether a new coverage map\index{coverage map} is interesting, the fuzzer compares it against
a global coverage map accumulated over the whole campaign.
$t^\prime$ is then deemed interesting if it contains coverage not yet contained in the global coverage map.
In practice, coverage maps are implemented as fixed-size arrays indexed by a hash of the coverage element.
This enables constant-time coverage lookups, but may cause coverage information to alias due to hash collisions~\cite{gan_collafl_2018}.

\subsubsection{Instrumentation}

Collecting coverage information requires instrumenting the SUT.
There are several categories of instrumentation in use~\cite{fioraldi_afl_nodate}:

\begin{enumerate}
  \item \emph{Compile-time instrumentation}\index{instrumentation!compile-time}
    inserts instrumentation during compile time of the SUT.
    This is very cheap and flexible, but requires access to the source code.

  \item \emph{Static binary rewriting}\index{instrumentation!static binary rewriting}
    injects the same code into an existing binary and removes the requirement
    of having source code available, but requires access to information that is not
    always available like libraries that are loaded dynamically.

  \item \emph{Dynamic binary translation}\index{instrumentation!dynamic binary translation}
    instead instruments blocks as they are translated at runtime, which works
    on arbitrary binaries and across architectures but costs a slowdown
    relative to native execution~\cite{bellard_qemu_2005}.

  \item \emph{Hardware tracing}\index{instrumentation!hardware tracing}
    makes use of built-in hardware instrumentation and tracing mechanisms like
    Intel Processor Trace~\cite{zhang_ptfuzz_2018}. This entails low overhead,
    but restricts instrumentation to host-native binaries.
\end{enumerate}

\subsection{AFL++}\label{sec:fuzzing-aflpp}

The most prominent instantiation of \cref{alg:cgf} is
AFL++\index{AFL++}~\cite{fioraldi_afl_nodate}, a reengineered fork of Zalewski's
now-unmaintained American Fuzzy Lop.
AFL++ consolidates various forks and research prototypes and exposes them as switchable modules~\cite{fioraldi_afl_nodate}.
It is therefore a customary baseline in the fuzzing literature.

\subsubsection{Test case selection}\label{sec:afl-test-case-selection}

AFL++ keeps the retained test cases in a queue and assigns to each a performance score.
It maintains a small set of \emph{favored} entries\index{favored entry} that together account for every edge observed so far.
\textsc{ChooseNext} is biased towards the favored entries~\cite{fioraldi_afl_nodate}.
A complementary \emph{trimming} stage\index{trimming} shrinks each queue entry to improve execution speed while preserving coverage.

For implementing \textsc{AssignEnergy}, AFL++ makes use of various power schedules
based on AFLFast~\cite{bohme_coverage-based_2016}. These include \texttt{fast},
\texttt{coe}, \texttt{explore}, \texttt{quad}, \texttt{lin} and
\texttt{exploit}, and are functions of per-seed statistics.
The default power schedule is \texttt{explore}.
It assigns a power proportional to the performance score $\alpha$:

\[
  p(t)=\frac{\alpha(t)}{\beta}\text{, where }\beta > 1
\]

On top of these, AFL++ adds \texttt{mmopt}, which raises the score of the newest queue
entries to encourage exploration, as well as
\texttt{rare}, which puts focus on queue entries that hit rare edges~\cite{fioraldi_afl_nodate,lemieux_fairfuzz_2018}.
It is noteworthy that a fixed schedule is picked before the start of the campaign, disallowing any adaptation to the target at runtime.

\subsubsection{Test case mutation}

AFL's mutation pipeline has two stages~\cite{fioraldi_afl_nodate}.
The \emph{deterministic} stage\index{mutation!deterministic stage} applies single, reproducible mutations systematically along
the input.
These include bit flips, arithmetic increments and decrements over byte, word and
double-word windows, and substitution of a fixed set of \enquote{interesting} constants
such as $0$, $-1$ or \texttt{INT\_MAX}.
The \emph{havoc} stage\index{mutation!havoc} stacks a random number of randomly chosen mutations, and unlike the
deterministic stage may also change the length of the input by inserting or deleting
byte blocks.
A later \emph{splicing} stage\index{mutation!splicing} merges two queue entries and applies havoc to the result.

AFL++ adds three further mutations that can be combined with the above.
\emph{MOpt}\index{MOpt}~\cite{lyu_mopt_nodate} replaces the fixed operator distribution with one
searched online by particle swarm optimization.
Another module consists of RedQueen's\index{RedQueen} \emph{input-to-state} (I2S)
correspondence\index{input-to-state correspondence}~\cite{aschermann_redqueen_2019}.
By tracking values used in comparison operations and comparing them with input values,
it can locate the responsible bytes in the input and substitute the value the comparison is testing for.
This allows overcoming magic byte and checksum checks without the overhead of concolic execution.

Another addition is a custom mutator API.
It allows external modules to take over parts of the loop through the \texttt{AFL\_CUSTOM\_MUTATOR\_LIBRARY} or \texttt{AFL\_PYTHON\_MODULE}
environment variables.
Rainfuzz, for instance, runs its policy-gradient-trained heat map over byte
positions alongside AFL++ through this interface~\cite{binosi_rainfuzz_2023}.

\subsubsection{Test case retention}

AFL implements \textsc{IsInteresting} through a hybrid metric combining edge coverage with
the number of times each edge was taken~\cite{fioraldi_afl_nodate}.
Basic blocks are assigned identifiers at instrumentation time.
Edges are represented by a hash of the identifiers of their source and destination blocks.
The instrumentation increments one byte per edge in a shared-memory bitmap that is reset before every execution.
After the run, each counter is quantized into one of eight power-of-two buckets and
compared against a bitmap accumulated over the campaign.
A test case is retained if it enters a new bucket for at least one edge.
The bitmap has a fixed size, so distinct edges may alias onto the same byte due to hash collisions.

The counter wraps around at 256 edge hits.
AFL++ offers a \emph{saturated counters} mode in which entries are frozen at 255,
as well as a \emph{NeverZero} mode\index{NeverZero} in which the carry flag is folded in so that
an overflowing counter never reaches zero.
NeverZero was shown to improve the visibility of hidden edges with a multiple of
256 hits, and is the default on most instrumentation backends~\cite{fioraldi_afl_nodate}.

\subsubsection{Instrumentation}

AFL++ supports six instrumentation backends: an LLVM pass, a GCC plugin, the legacy
\texttt{afl-gcc} wrapper, QEMU~\cite{bellard_qemu_2005}, \emph{unicornafl}, and QBDI.
This covers compile-time instrumentation (LLVM, GCC) and dynamic binary
translation (QEMU for user-space binaries, unicornafl for raw firmware blobs),
with QBDI adding dynamic binary instrumentation for closed-source Android
native libraries.
Some AFL++ components are backend-specific: CmpLog, for instance, is available only under LLVM and QEMU.
The choice of instrumentation therefore constrains which other components can be used.

Granularity is configurable between block, edge and path level.
\emph{Context-sensitive} edge coverage\index{coverage!context-sensitive} XORs a block's identifier with an identifier of the
current call context, so that the same edge reached from different callers is counted
separately, improving edge discrimination and reducing aliasing.
\emph{Ngram} coverage\index{coverage!ngram} hashes the destination block together with the $N-1$ preceding blocks
for $N$ between 2 and 16, approximating path coverage.
\emph{InsTrim} adds instrumentation only to blocks that have been deemed interesting
by a prior static analysis pass, reducing instrumentation overhead~\cite{hsu2018instrim}.

\subsection{Strengths and Limitations of Coverage-Guided Fuzzing}\label{sec:fuzzing-limitations}

\subsubsection{Strengths}

Collecting coverage information can usually be done using lightweight instrumentation.
This makes coverage-guided fuzzing extremely fast, especially in comparison to
methods using some form of symbolic execution and constraint solving~\cite{godefroid_sage_2012,manes_art_2019}.
Since execution throughput is usually correlated with the effectiveness of a fuzzer,
this is an important advantage~\cite{fioraldi_afl_nodate,manes_art_2019}.
It also does not require an input model: an iteratively refined corpus can
implicitly generate inputs that adhere to the specification and grammar of the SUT.
This makes CGF applicable to a wide range of targets~\cite{manes_art_2019}.
It works well at scale and has found tens of thousands of
bugs~\cite{arya_oss-fuzz_nodate} in a wide range of software.

\subsubsection{Limitations}

\textcite{bohme_fuzzing_2020} have shown that although linearly increasing the fuzzing budget
produces a linear reduction in time required to find a given set of bugs, actually
discovering new bugs requires an exponentially increasing fuzzing budget.
AFLFast attributes this plateau to the budget being spent on high-frequency
paths~\cite{bohme_coverage-based_2016}.

CGF is also blind to any input-coverage correspondence~\cite{manes_art_2019}, therefore
discarding a large amount of potentially useful information.
Gradient-guided optimization techniques have been shown to outperform evolutionary
techniques like the ones used by AFL.
\textcite{she_neuzz_2019,she_mtfuzz_2020} show that learning an approximation to a SUT's branching
behavior improves fuzzing performance.
\textcite{gan_greyone_nodate} use taint analysis to discover more edges than baseline AFL.
Further approaches are surveyed by~\textcite{wang_systematic_2020} and~\textcite{qiu_survey_2025}.

Magic bytes and checksums present a major barrier to the success of CGF since
they are highly unlikely to be found by random mutation (\cref{sec:fuzzing-input-generation}).
Overcoming these checks can be done by augmenting CGF with symbolic
execution~\cite{stephens_driller_2016} or by learned input-coverage
correspondence using the aforementioned methods~\cite{aschermann_redqueen_2019}.

Coverage itself might also be an insufficient objective to optimize for. While
it has been empirically proven to be effective for general SUTs, in some use
cases like embedded systems, corruption of a SUT's state is not related to its
coverage~\cite{muench_what_2018} or might require additional sanitization
efforts~\cite{zhang_debloating_nodate}.

Coverage is also a lossy representation that only approximates the SUT's state
at a certain point and might suffer from aliasing due to hash collisions~\cite{gan_collafl_2018}.

Fuzzing itself can be modeled as a decision problem under uncertainty whose
solution is strongly target-dependent~\cite{woo_scheduling_2013,bohme_coverage-based_2016}.
Hand-tuned and precomputed seed, power and mutation schedules therefore cannot be optimal.
MOpt~\cite{lyu_mopt_nodate} is included in AFL++ and aims to optimize the distribution of
mutation operators using particle swarm optimization, but is still restricted to
a fixed set of available mutation operators.
AFL++'s \textsc{ChooseNext} and \textsc{AssignEnergy} are adaptive, but depend on a few hand-crafted metrics,
none of which uses raw coverage information (\cref{sec:afl-test-case-selection}).

\subsection{Stateful and Protocol Fuzzing}\label{sec:fuzzing-stateful}

CGF models the SUT as a pure function mapping an input to a coverage.
A new input requires the SUT to be reset to its initial state.
For stateful SUTs, the coverage depends on the current input as well as the
history of previous inputs.
Stateful targets\index{fuzzing!stateful} are especially prevalent in IoT applications and present a
significant attack surface~\cite{zhang_iot_2014,sadeghi_security_2015}.
Examples include industrial control systems~\cite{borcherding_bandits_2023},
network protocols like QUIC~\cite{ang_quic-fuzz_2025} and OPC~UA,
operating system kernels~\cite{wang_syzvegas_2021}, and web applications~\cite{doupe_enemy_2012}.

We distinguish three kinds of state:

\paragraph{Fuzzer state}
The current fuzzing configuration and corpus.\index{state!fuzzer}

\paragraph{Program state}
The state of the SUT as seen from the execution engine.\index{state!program} It is made up
of memory and register file state as well as peripheral states like filesystems or network interfaces.

\paragraph{Protocol state}
The state inside the protocol as given by a protocol specification.\index{state!protocol} For a SUT initiating
a TCP connection, this might include the protocol state after a \texttt{SYN} has been sent.

\subsubsection{Challenges}

Fuzzers like AFL++ were mainly designed to test stateless programs and
therefore cannot deal natively with highly stateful SUTs.
Coverage in this case is not sufficient to identify the state the program is currently
in since it does not contain enough information about the input history.
In addition, protocol state is not observable and must be inferred, for instance from response codes,
which introduces additional uncertainty.
Further, most execution engines only support execution from the beginning, while
a stateful fuzzer would ideally have access to an arbitrary snapshot-and-reset
mechanism.
Resetting the environment can be difficult, since targets like FTP servers
might create lots of files that then need to be removed again.
We provide such an execution environment in~\cref{sec:impl-platform}.

AFLNet\index{AFLNet}~\cite{pham_aflnet_2020} is the canonical greybox fuzzer for network protocols\index{fuzzing!protocol}
and works by combining traditional CGF with automatic state inference to guide exploration
towards more important code paths.
Various prior works describe how a protocol state machine can be learned~\cite{dongol_grey-box_2020,marksteiner_learning_2025,giraud_l-scale_2026}.
A detailed survey of fuzzers that incorporate protocol state information is given by~\textcite{daniele_fuzzers_2023}.
